% !TEX root = ../ComputationalOTFnT.tex

\chapter{アルゴリズムの基礎}
\label{c-algo-basics} % a label for the chapter, to refer to it later

本章では、主問題~\eqref{eq-mk-discr}またはその双対~\eqref{eq-dual}で記述される最適輸送の離散定式化を解くために用いられる、組合せ最適化および線形計画法における最も一般的なアルゴリズム的手法について述べる。

これらのアルゴリズムの起源は第二次世界大戦にまで遡ることができ、戦前の\citeauthor{tolstoi1930methods}の先駆的な研究 \citeyearpar{tolstoi1930methods}、あるいは戦中に~\citet{Hitchcock41}および\citet{Kantorovich42}が、利用可能な資源を消費地へ最適に配分するという一般的な問題を定式化した時期に始まる。これらの定式化、および後の~\citet{koopmans1949optimum}による貢献は、いずれもその問題を解くための\emph{証明可能に}正しいアルゴリズムを提供するには至らなかった（巡回違反法は~\citet{tolstoi1939metody}によってヒューリスティックとして既に提案されていた）。線形計画法の分野が十分に成熟し、シンプレックス法が提案されるまで、これらの問題を厳密に解くことはできなかった。

線形計画法の目的は、目的関数が線形であり、制約が関心のある変数に関する線形（不）等式であるような最適化問題を解くことである。最適輸送問題はこの記述に適合するため、この広いクラスの問題の特殊な場合である。しかしながら、最適輸送はすべての線形計画の中でも真に特別であると言える。第一に、\citeauthor{dantzig1991}が線形計画を解こうとした初期の動機は、輸送問題を解くこととと密接に関連していた~\citep[p. 210]{dantzig49econometrica}。第二に、特殊な場合に過ぎないにもかかわらず、最適輸送問題は最適化の分野で注目され続けた。それは、最適輸送問題が最小費用ネットワークフローとして知られる線形計画の重要なクラスと関連し、実際には等価であることが、~\citet{ford1962flows}の結果により~\citep[p. 213, Lem. 9.3]{korte2012combinatorial}間もなく理解されたためである。そのため、OT問題は数理計画法の誕生~\citep{Dantzig51}以来、特に注目されてきており、新しい読者に最適化を紹介するために今なお広く用いられている~\citep[\S1, p. 4]{nocedal}。

\section{Kantorovichの線形計画}\label{s-kantorovich}
式~\eqref{eq-mk-discr}で主OT問題を既に導入した：
\eql{\label{eq-mk-discr-algo}
	\MKD_{\C}(\a,\b) =
	\umin{\P \in \CouplingsD(\a,\b)}
		\sum_{i\in \range{n}, j \in \range{m}} \C_{i,j} \P_{i,j}.
}
線形計画法の文献との関連を明確にするために、上記の式を\emph{標準}形の線形計画、すなわち線形目的関数、行列と定数ベクトルで定義される等式制約、および変数の非負制約を持つ線形計画として定式化できる。$\Identity_n$をサイズ$n$の単位行列、$\otimes$をクロネッカー積とする。${(n+m) \times nm}$行列
$$
\mathbf{A}= \begin{bmatrix}
	\transp{\ones_{n}} \otimes \Identity_m \\
	\Identity_n \otimes \transp{\ones_{m}}
\end{bmatrix}\in\RR^{(n+m)\times nm}
$$
は、任意の$\P$が$\CouplingsD(\a,\b)$に属するために満たすべき行和および列和の制約を符号化するために用いることができる。そのために、行列$\P\in\RR^{n\times m}$をベクトル$\p\in\RR^{nm}$として表し、$\p$の第$i+n(j-1)$要素が$\P_{ij}$に等しくなるようにする（$\P$は列方向に列挙される）と、以下の同値が得られる：
$$\P\in \RR^{n\times m}\in \CouplingsD(\a,\b) \Leftrightarrow \p\in\RR^{nm}_+, \mathbf{A}\p = \bigl[\begin{smallmatrix}\a\\ \b \end{smallmatrix} \bigr].$$
したがって、元の最適輸送問題を以下のように書くことができる：
\eql{\label{eq-mk-discr-primal}
	\MKD_{\C}(\a,\b) =
	\umin{\substack{\p \in \RR^{nm}_+\\
		\mathbf{A}\p = \bigl[\begin{smallmatrix}\a\\ \b \end{smallmatrix} \bigr]}} \transp{\cc}\p,
}
ここで$nm$次元ベクトル$\cc$はコスト行列$\C$の列を積み重ねたものに等しい。

\begin{rem}\label{rem-transportation-polytope} 上記の$n+m$個の制約のうち一つは冗長であること、言い換えれば、行列$A$の行ベクトルは線形独立ではないことに注意されたい。実際、最初の$n$行と後続の$m$行をそれぞれ足し合わせると同じベクトルが得られる（すなわち$A \bigl[\begin{smallmatrix} \ones_n\\\zeros_m \end{smallmatrix}\bigr] = A \bigl[\begin{smallmatrix} \zeros_n\\\ones_m\end{smallmatrix}\bigr]=\transp{\ones_{nm}}$）。$A$の一つの行とそれに対応する$\bigl[\begin{smallmatrix}\a\\ \b \end{smallmatrix}\bigr]$のエントリを除去しても、適切に定義された線形系が得られることが示せる。簡潔さのため、また$\a$と$\b$を非対称に扱うことを避けるため、以下では冗長な定式化を維持し、いくつかの計算で退化が生じることに留意する。\end{rem}

式~\eqref{eq-mk-discr-primal}に対応する双対問題は、線形計画法における双対性~\citep[p. 143]{bertsimas1997introduction}に従い、以下のように定義される：
\eql{\label{eq-mk-discr-dual}
	\MKD_{\C}(\a,\b) =
	\umax{\substack{\hD \in \RR^{n+m}\\
		\transp{A}\hD \leq \cc}}  \transp{\bigl[\begin{smallmatrix}\a\\ \b \end{smallmatrix} \bigr]}\hD.
}

この計画問題は式~\eqref{prop-duality-discr}で示されたものと正確に等価であることに注意されたい。


\begin{rem}\label{rem-duality} 上記の双対性の結果の簡潔な導出を示す。これは注~\ref{rem-kantorovich-dual}で展開された議論の直接的な定式化と見なすことができる。強双対性、すなわち主問題~\eqref{eq-mk-discr-primal}と双対問題~\eqref{eq-mk-discr-dual}の最適値が実際に一致するという事実は、より長い証明を要する~\citep[\S4.10]{bertsimas1997introduction}。記法を簡略化するため、$\q=\bigl[\begin{smallmatrix}\a\\ \b \end{smallmatrix} \bigr]$と書く。
ここで、最適輸送問題の緩和された主問題を考える。制約$\mathbf{A}\p=\q$はもはや必ずしも強制されず、代わりに任意のコストベクトル$\hD\in\RR^{n+m}$によってパラメータ化されたコスト$\transp{\hD}(\mathbf{A}\p-\q)$が課される。
この緩和の最適値はコストベクトル$\hD$に直接依存し、以下のように書ける：
$$\LagrangeMKD(\hD)\eqdef \umin{\p \in \RR^{nm}_+} \transp{\cc}\p-\transp{\hD}(\mathbf{A}\p-\q).$$
まず、この緩和問題は$\p$に対する周辺制約を持たないことに注意する。この最小化はより多くの$\p$の解を許容するため、$\LagrangeMKD(\hD)$は$\bar{z}=\MKD_{\C}(\a,\b)$よりも小さくなることが期待される。実際、主問題\eqref{eq-mk-discr-algo}の任意の最適解を$\p^\star$と書くと、
$$\umin{\p \in \RR^{nm}_+} \transp{\cc}\p-\transp{\hD}(\mathbf{A}\p-\q) \leq \transp{\cc}\p^\star-\transp{\hD}(\mathbf{A}\p^\star-\q)=\transp{\cc}\p^\star=\bar{z}.$$
したがって、上記のアプローチは、任意のコストベクトル$\hD$に対して元の問題~\eqref{eq-mk-discr-algo}の最適な上界を計算するための問題を定義する。この関数は$\MKD$のLagrange双対関数と呼ばれる。
双対理論の目的は、\emph{任意の}コストベクトル$\hD$に対して$\LagrangeMKD$を\emph{最大化}することにより、最良の下界$\underline{z}$を計算することである。すなわち
$$
	\underline{z}=\umax{\hD} \pa{ \LagrangeMKD(\hD)=\umax{\hD} \transp{\hD}\q + \umin{\p \in \RR^{nm}_+} \transp{(\cc-\transp{A}\hD)}\p }.
$$
$\p$に関する最小化を含む第二項は、$\transp{\cc}- \transp{A}\hD$のいずれかの座標が負であれば$-\infty$になることが容易に示される。
実際、例えばある添字$i\leq n+m$に対して$\cc_i-(\transp{A}\hD)_i<0$であれば、$\p$として正準基底ベクトル$\e_i$に任意の大きな正の値を掛けたものを取ることで、非有界な値が得られる。
下界$\LagrangeMKD(\hD)$を最大化しようとする場合、$\transp{A}\hD\leq \cc$を満たすベクトル$\hD$に制限することが合理的であり、その場合の最良の下界は$$\underline{z}=	\umax{\substack{\hD \in \RR^{n+m}\\\transp{A}\hD \leq \cc}}\transp{\hD}\q.$$となる。これにより弱双対性の結果、すなわち$\underline{z}\leq \bar{z}$が証明された。
\end{rem}


\section{$\C$変換}\label{sec-c-transforms} 本節では、双対最適輸送問題~\eqref{eq-mk-discr-dual}の重要な性質を示す。この性質は\S\ref{s-c-transform}における半離散最適輸送問題で用いる際により重要な意味を持つ。本節は、双対変数を行和と列和の制約に従って分離する元の定式化~\eqref{eq-dual}に基づいている：
\eql{\label{eq-dual-discrete-split}\MKD_\C(\a,\b) = \umax{(\fD,\gD) \in \PotentialsD(\C)} \dotp{\fD}{\a} + \dotp{\gD}{\b}.
}
任意の双対実行可能対$(\fD,\gD)$を考える。$\fD$の値を「固定」すれば、$\gD$に対する最良のベクトル解は$\fD$の$\C$変換ベクトルであることに気付く。これを$\fD^{\,\C}\in\RR^m$と書き、以下で定義される：
\eq{(\fD^{\,\C})_j = \min_{i\in \range{n}} \C_{ij}-\fD_i,}
実際、$(\fD,\fD^{\,\C})\in\PotentialsD(\C)$であること、および$\fD^{\,\C}$がこの制約を満たす最大のベクトルであることが容易に証明できる。したがって、
\eq{\dotp{\fD}{\a} + \dotp{\gD}{\b} \leq \dotp{\fD}{\a} + \dotp{\fD^{\,\C}}{\b}.}
が成り立つ。この結果により、まず双対問題を単一変数$\fD$で表される区分的アフィン凹最大化問題として再定式化できる：
\eql{\label{eq-semidual-discret}\MKD_\C(\a,\b) = \umax{\fD\in\RR^{n}} \dotp{\fD}{\a} + \dotp{\fD^{\,\C}}{\b}.}

この結果を一旦脇に置くと、$\gD$の値を「固定」して$\gD$の$\bar{\C}$変換、すなわちベクトル$\gD^{\bar{\C}}\in\RR^n$を以下で定義すれば、もちろん同じ議論が成り立つ：
\eq{(\gD^{\bar{\C}})_i = \min_{j\in \range{m}} \C_{ij}-\gD_j,}
目的関数の増加は異なり、
\eq{\dotp{\fD}{\a} + \dotp{\gD}{\b} \leq \dotp{\gD^{\bar{\C}}}{\a} + \dotp{\gD}{\b}.}
与えられた$\fD$から出発して、$\fD$を改善するために$\C$変換と$\bar{\C}$変換を複数回交互に適用したくなる。実際、以下の不等式の列が得られる：
\eq{
	\dotp{\fD}{\a} + \dotp{\fD^{\,\C}}{\b} \leq \dotp{\fD^{\,\C\bar{\C}}}{\a} + \dotp{\fD^{\,\C}}{\b}\leq  \dotp{\fD^{\,\C\bar{\C}}}{\a} + \dotp{\fD^{\,\C\bar{\C}\C}}{\b}\leq  \dots
}
これらの反復の各ステップで目的関数が真に増加することを期待するかもしれない。しかし、$\C$変換と$\bar{\C}$変換を交互に適用するとすぐに停滞するため、これはうまくいかない。

\begin{prop}\label{prop-ccc-2}以下の等式が成り立つ。ここでベクトル間の不等号は要素ごとに理解されるものとする：
	\begin{enumerate}[label=(\roman*)]
	\item $\fD\leq \fD\,' \Rightarrow \fD^{\,\C}\geq \fD\,'^{\,\C}$,
	\item $\fD^{\,\C\bar{\C}} \geq \fD$, $\gD^{\bar{\C}\C} \geq \gD$,
	\item $\fD^{\,\C\bar{\C}\C}=\fD^{\,\C}.$
\end{enumerate}
\end{prop}
\begin{proof} 第一の不等式は$\C$変換の定義から従う。$\fD^{\,\C\bar{\C}}$の定義を展開すると
\eq{\left(\fD^{\,\C\bar{\C}}\right)_i= \min_{j\in \range{m}} \C_{ij}-\fD^{\,\C}_j = \min_{j\in \range{m}} \C_{ij}- \min_{i'\in \range{n}} \C_{i'j}-\fD_{i
'}.}
ここで、$-\min_{i'\in \range{n}} \C_{i'j}-\fD_{i'} \geq -(\C_{ij}-\fD_i)$であるから、
\eq{\left(\fD^{\,\C\bar{\C}}\right)_i \geq \min_{j\in \range{m}} \C_{ij}- \C_{ij}+\fD_i = \fD_i.}
が回復される。関係式$\gD^{\bar{\C}\C} \geq \gD$も同様に得られる。次に、$\gD=\fD^{\,\C}$とおく。すると、$\gD^{\bar{\C}}=\fD^{\,\C\bar{\C}}\geq \fD$である。したがって、結果(i)を用いると$\fD^{\,\C\bar{\C}\C}\leq \fD^{\,\C}$が得られる。結果(ii)により$\fD^{\,\C\bar{\C}\C}\geq \fD^{\,\C}$であるから、等式が証明される。
\end{proof}
%Applying the $\C\bar{\C}$ transform more than once has therefore no impact.


\section{相補スラック性}\label{s-complementary}

主問題~\eqref{eq-mk-discr-primal}と双対問題~\eqref{eq-mk-discr-dual},~\eqref{eq-dual}は独立に解くことができ、最適な主変数$\P^{\star}$および双対変数$(\fD^{\star},\gD^{\star})$が得られる。以下の命題はそれらの関係を特徴づける。

\begin{prop}\label{prop-primal-dual-optimal}$\P^\star$および$\fD^\star,\gD^\star$をそれぞれ主問題~\eqref{eq-dual-generic}および双対問題~\eqref{eq-mk-discr}の最適解とする。このとき、任意の対$(i,j)\in\range{n}\times\range{m}$に対して、$\P^\star_{i,j}(\C_{i,j}-\fD^\star_i+\gD^\star_j)=0$が成り立つ。言い換えれば、$\P^\star_{i,j}>0$ならば必然的に$\fD^\star_i+\gD^\star_j=\C_{i,j}$であり、$\fD^\star_i+\gD^\star_j<\C_{i,j}$ならば必然的に$\P^\star_{i,j}=0$である。
\end{prop}
\begin{proof} 強双対性により$\dotp{\P^\star}{\C} = \dotp{\fD^\star}{\a}+\dotp{\gD^\star}{\b}$が成り立つ。$\P^\star\ones_m = \a$および$\transp{\P^\star}\ones_n = \b$であることを想起すると、
	$$\begin{aligned}\dotp{\fD^\star}{\a}+\dotp{\gD^\star}{\b} &= \dotp{\fD^\star}{\P^\star\ones_m}+\dotp{\gD^\star}{\transp{\P^\star}\ones_n}\\&=  \dotp{\fD^\star\transp{\ones_m}}{\P^\star}+\dotp{\ones_n\transp{\gD^\star}}{\P^\star},\end{aligned}$$
	となり、$$\dotp{\P^\star}{\C- \fD^\star\oplus \gD^\star}=0.$$
が得られる。$(\fD^\star,\gD^\star)$は双対制約の多面体~\eqref{eq-feasible-potential}に属するため、行列$\C- \fD^\star\oplus \gD^\star$の各要素は必然的に非負である。したがって、$\P$のすべての要素が非負であることから、上記の内積がゼロに等しいという制約は、$\P_{i,j}>0$である任意の添字対$(i,j)$に対して$\C_{i,j}-(\fD_i+\gD_j)$がゼロでなければならないこと、および$\C_{i,j}>\fD_i+\gD_j$である任意の添字対$(i,j)$に対して$\P_{i,j}=0$でなければならないことを強制する。
\end{proof}

逆の結果もまた成り立つ。まず、主問題と双対問題の二つの変数が相補的であるという概念を定義する。

\begin{defn}\label{def-complementary}
行列$\P\in\RR^{n\times m}$とベクトルの対$(\fD,\gD)$が$\C$に関して相補的であるとは、$\P_{i,j}>0$であるすべての添字対$(i,j)$に対して$\C_{i,j}=\fD_i+\gD_j$も成り立つことをいう。
\end{defn}

主変数と双対変数の実行可能な対が相補的であれば、それらが最適であると結論付けることができる。

\begin{prop}\label{prop-primal-dual-optimality} $\P$と$(\fD,\gD)$がそれぞれ主問題~\eqref{eq-dual-generic}と双対問題~\eqref{eq-mk-discr}の相補的かつ実行可能な解であるならば、$\P$と$(\fD,\gD)$はともに主最適かつ双対最適である。
\end{prop}
\begin{proof}
	弱双対性により、
$$\MKD_{\C}(\a,\b)\leq \dotp{\P}{\C}=\dotp{\P}{\fD\oplus\gD}=\dotp{\a}{\fD}+\dotp{\b}{\gD}\leq \MKD_{\C}(\a,\b)$$
が成り立ち、したがって$\P$と$(\fD,\gD)$はそれぞれ主最適および双対最適である。
\end{proof}


\todoK{
\begin{prop}[主双対関係]
	\todo{something on this?}
\end{prop}
}

\section{輸送多面体の頂点}\label{s-extremal}

凸集合の頂点または端点とは、形式的にはその集合に属する点$\mathbf{x}$であって、その集合に属する$\mathbf{y}$と$\mathbf{z}$が存在し$\mathbf{x}=(\mathbf{y}+\mathbf{z})/2$が成り立つならば、必然的に$\mathbf{x}=\mathbf{y}=\mathbf{z}$であるような点をいう。空でなく有界な実行可能集合を持つ線形計画は、その実行可能集合の頂点（または端点）で最小値を達成する~\citep[p.~65, Theo.~2.7]{bertsimas1997introduction}。主最適輸送問題~\eqref{eq-mk-discr-primal}の実行可能集合$\CouplingsD(\a,\b)$は有界であるため、最適な$\P$の探索を多面体$\CouplingsD(\a,\b)$の端点の集合に限定できる。$\CouplingsD(\a,\b)$で端点となる行列$\P$は、広範な研究の対象となってきた興味深い構造を持つ~\citep[\S8]{brualdi2006combinatorial}。その構造を記述するには、輸送問題を二部グラフの形式で表す必要がある。

\subsection{$\CouplingsD(\a,\b)$のすべての頂点の台の木構造}
$V=(1,2,\dots,n)$と$V'=(1',2',\dots,m')$を二つのノード集合とする。集合$V'$のラベルには$V$のラベルと区別するためにプライムを付ける。それらの和集合$V\cup V'$は$n+m$個のノードを持ち、それらの間のすべての$nm$本の有向辺の集合$\alledges$を$\{ (i,j'), i \in \range{n}, j\in \range{m}\}$とする（ここでは整数$j\leq m$にプライムを付けて$V'$の$j'$を構成する）。各辺$(i,j')$に対して対応するコスト値$\C_{ij}$を関連付ける。$V$と$V'$の間の完全二部グラフ$\mathcal{G}$は$(V\cup V',E)$である。輸送計画はこのグラフ上の、ソース（各ノード$i$から$\a_i$が流出する）およびシンク（各ノード$j'$に$\b_j$が流入する）制約を満たすフローであり、図~\ref{fig-simplex}に非形式的に描かれている。$\CouplingsD(\a,\b)$の端点は以下の性質を持つ~\citep[p.~338,~Theo.~8.1.2]{brualdi2006combinatorial}。

\begin{figure}[h!]
\centering
\includegraphics[width=\linewidth]{algo-basics-tree-solution/networksol}
\caption{\label{fig-simplex}
二部ネットワークフロー問題としての最適輸送問題。ここで$n=3,m=4$である。ソースヒストグラム$\a$のすべての座標は左側のソースノード$1,2,3$として描かれ、ターゲットヒストグラム$\b$のすべての座標はノード$1',2',3',4'$として描かれている。このグラフはすべてのソースノードがすべてのターゲットノードに接続され、追加の辺がないという意味で二部である。各辺$(i,j')$にはコスト$\C_{ij}$が関連付けられている。実行可能なフローが右側に表示されている。命題~\ref{prop-extremal}は、少なくとも一つの巡回$((1,1'),(2,1'),(2,4'),(1,4'))$が存在するため、このフローが端点ではないことを示している。
}
\end{figure}


\begin{figure}[h!]
\centering
\includegraphics[width=\linewidth]{algo-basics-tree-solution/tree-sol2}
\caption{\label{fig-perturb}
台のグラフに巡回を持つ解$\P$は、$\P$がそれらの平均となるような二つの実行可能解$\Q$および$\RRd$を得るために摂動でき、それにより$\P$が端点であることが否定される。
}
\end{figure}

\begin{prop}[端点解]\label{prop-extremal}
$\P$を多面体$\CouplingsD(\a,\b)$の端点とする。$S(\P)\subset \alledges$を辺の部分集合$\{(i,j'), i\in\range{n}, j\in\range{m} \text{ such that } \P_{ij}>0\}$とする。このとき、グラフ$G(\P)\eqdef (V\cup V',S(\P))$は巡回を持たない。特に、$\P$は$n+m-1$個を超える非零要素を持つことができない。
 \end{prop}
\begin{proof}
背理法による。$\P$が多面体$\CouplingsD(\a,\b)$の端点であり、その対応する辺の集合$S(\P)$（簡潔に$F$と書く）に対してグラフ$G=(V\cup V',F)$が巡回を含むと仮定する。すなわち、$k>1$と相異なる添字の列$i_1,\dots,i_{k-1}\in\range{n}$および$j_1,\dots,j_{k-1}\in\range{m}$が存在し、以下に示す辺の集合$H$が$F$の部分集合を形成する。$$H=\left\{(i_1,j_1'), (i_2,j_1'), (i_2,j_2'),\dots,(i_k,j_k'),(i_1,j_k')\right\}.$$
	ここで、$\P=(\Q+\RRd)/2$を満たす二つの実行可能行列$\Q$と$\RRd$を構成する。そのために、$H$に対応する\emph{有向}巡回$\bar{H}$、すなわち対の列$i_1\rightarrow j_1', j_1' \rightarrow i_2, i_2 \rightarrow j_2',\dots, i_k \rightarrow j_k' , j_k' \rightarrow i_1$、および基本フロー量$\varepsilon < \min_{(i,j')\in F}\P_{ij}$を考える。
	摂動行列$\E$を、その$(i,j)$要素が$i\rightarrow j' \in \bar{H}$ならば$\varepsilon$、$j\rightarrow i'\in \bar{H}$ならば$-\varepsilon$、それ以外は零であるように定義する。行列$\Q=\P+\E$と$\RRd=\P-\E$を図~\ref{fig-perturb}に示されるように定義する。$\varepsilon$が十分小さいため、$\Q$と$\RRd$のすべての要素は非負である。構成上、$\E$はすべての要素がゼロの行（もしくは列）を持つか、$i_1,\dots,i_k$（もしくは$j_1,\dots,j_k$）で添字付けられた行（もしくは列）について、一つの要素が$\varepsilon$でもう一つが$-\varepsilon$であるかのいずれかである。したがって、$\E$は$\E\ones_{m}=\zeros_{n}$かつ$\transp{\E}\ones_{n}=\zeros_{m}$を満たし、$\Q$と$\RRd$は$\P$と同じ周辺分布を持ち、よって実行可能である。最後に$\P=(\Q+\RRd)/2$であり、$\Q,\RRd\ne \P$であるから、$\P$が端点であるという事実に矛盾する。$k$個のノードを持ち巡回のないグラフは$k-1$本を超える辺を持てないため、$S(\P)$は$n+m-1$本を超える辺を持てず、したがって$\P$は$n+m-1$個を超える非零要素を持てない。
\end{proof}

\subsection{北西隅法}\label{subsec-northwest}
北西（NW）隅法は、最大$n+m$回の操作で多面体$\CouplingsD(\a,\b)$の頂点を生成するヒューリスティックである。このヒューリスティックは、次節で概説するネットワークシンプレックスのような、主変数に基づいて動作するアルゴリズムの初期化において役割を果たす。

この規則は、$\P_{1,1}$に可能な限り大きな値を与えることから始まり、$\min(\a_1,\b_1)$に設定する。各ステップにおいて、要素$\P_{i,j}$は行制約$i$、列制約$j$、または可能であれば両方を飽和させるように選ばれる。添字$i,j$は次のように更新される：第一の場合は$i$を増やし、第二の場合は$j$を増やし、第三の場合は$i$と$j$の両方を増やす。この規則は$\P_{n,m}$に値が割り当てられるまで続く。

形式的には、アルゴリズムは以下のように動作する：$i$と$j$を$1$に初期化し、$r\leftarrow\a_1,c\leftarrow \b_1$とする。$i\leq n$かつ$j\leq m$の間、$t\leftarrow \min(r,c)$、$\P_{i,j}\leftarrow t$、$r\leftarrow r-t$、$c\leftarrow s-t$とし、$r=0$ならば$i$を増やし、$i\leq n$ならば$r\leftarrow\a_i$と更新する；$c=0$ならば$j$を増やし、$j\leq n$ならば$c\leftarrow \b_j$と更新する；繰り返す。以下は$\a=[0.2,0.5,0.3]$および$\b=[0.5,0.1,0.4]$を仮定した場合のこの手順の例である：
$$\begin{aligned}\begin{bmatrix} \bullet & 0 & 0 \\ 0 & 0 & 0 \\ 0& 0 & 0\end{bmatrix} &\rightarrow \begin{bmatrix} 0.2 & 0 & 0 \\ \bullet & 0 & 0 \\ 0& 0 & 0\end{bmatrix} &\rightarrow \begin{bmatrix} 0.2 & 0 & 0 \\ 0.3 & \bullet & 0 \\ 0& 0 & 0\end{bmatrix}\\ &\rightarrow \begin{bmatrix} 0.2 & 0 & 0 \\ 0.3 &0.1 &\bullet \\ 0& 0 & 0\end{bmatrix} &\rightarrow \begin{bmatrix} 0.2 & 0 & 0 \\ 0.3 &0.1 &0.1 \\ 0& 0 & \bullet\end{bmatrix} & \rightarrow \begin{bmatrix} 0.2 & 0 & 0 \\ 0.3 &0.1 &0.1 \\ 0& 0 & 0.3\end{bmatrix}\end{aligned}$$
このヒューリスティックにより得られる一意な計画を$\NW(\a,\b)$と書く。

しかし、$\a$と$\b$の順序を任意に置換し、対応するNW隅テーブルを計算し、列と行の順序を再度反転してテーブルを$\CouplingsD(\a,\b)$に復元することで、はるかに多くのNW隅解を得られることに注意されたい：$\sigma=(3,1,2),\sigma'=(3,2,1)$とすると、$\a_\sigma=[0.3,0.2,0.5], \b_{\sigma'}=[0.4,0.1,0.5]$、および$\sigma^{-1}=(2,3,1),\sigma'=(3,2,1)$となる。以下を観察されたい：
\begin{gather*}
\NW(\a_\sigma,\b_{\sigma'}) = \begin{bmatrix} 0.3 & 0 & 0 \\ 0.1 & 0.1 & 0 \\ 0& 0 & 0.5\end{bmatrix} \in \CouplingsD(\a_\sigma,\b_{\sigma'}),\\
\NW_{\sigma^{-1}\sigma'^{-1}}(\a_\sigma,\b_{\sigma'})= \begin{bmatrix} 0 & 0.1 & 0.1 \\ 0.5 & 0 & 0 \\ 0& 0 & 0.3\end{bmatrix}\in \CouplingsD(\a,\b).
\end{gather*}

$\mathcal{N}(\a,\b)$をこの方法で生成可能なすべてのNW隅解の集合とする：
$$\mathcal{N}(\a,\b)\eqdef\{ \NW_{\sigma^{-1}\sigma'^{-1}}(r_\sigma,c_{\sigma'}), \sigma,\sigma'\in S_d\}.$$
すべてのNW隅解は、構成上、最大$n+m-1$個の非零要素を持つ。NW隅法は$\a_\sigma$と$\b_\sigma'$に対して構成上一意なテーブルを生成するが、同じテーブルを与える行・列置換の対$(\sigma,\sigma')$は指数的に多く存在し得る~\citep[p. 2]{stougie2002polynomial}。$\mathcal{N}(\a,\b)$は$\CouplingsD(\a,\b)$の端点集合の（通常は真の）部分集合を形成する~\citep[Cor. 8.1.4]{brualdi2006combinatorial}。



\section{ネットワークシンプレックスのヒューリスティックな説明}\label{s-networksimplex}

グラフ$G(\P)=(V\cup V',S(\P))$が巡回を持たない実行可能行列$\P$を考える。$\P$は$n+m-1$個以下の非零要素を持ち、命題~\ref{prop-extremal}により$\CouplingsD(\a,\b)$の頂点である。
%
命題~\ref{prop-primal-dual-optimality}に従えば、$\P$が最適であることを証明するには、実行可能（すなわち$\C-\fD\oplus\gD$が非負の要素を持つ）かつ$\P$に相補的（$S(\P)$の添字対$(i,j')$が$\C_{i,j}=\fD_i+\gD_j$を満たす）な双対解$(\fD,\gD)$を得れば十分である。
%
ネットワークシンプレックスは二つの単純な原理に基づいている：各実行可能な主解$\P$に対して相補的な対$(\fD,\gD)$を関連付けることができる。
%
その対が実行可能であれば最適に到達している。そうでなければ、$\P$を実行可能性を保ちつつ修正し、その相補的な対$(\fD,\gD)$が実行可能性に近づくように変更することを考える。

\subsection{$\P$に相補的な双対対の取得}\label{subsec-obtaining}
シンプレックス法はまず、任意の端点解$\P$に相補的な双対変数の対$(\fD,\gD)$を関連付けることから始める。これは、$S(\P)$の任意の$(i,j')$に対して$\fD_i+\gD_j$が$\C_{i,j}$に等しくなるような二つのベクトル$\fD$と$\gD$を見つけることで単純に実行される。ただし、これ自体は$(\fD,\gD)$の実行可能性を保証しないことに注意する。

\begin{figure}[h!]
	\centering
	\includegraphics[width=\linewidth]{algo-basics-tree-solution/f_of_P}
	\caption{\label{fig-f_of_P}
	実行可能な輸送$\P$と、それに対応する辺の集合$S(\P)$およびグラフ$G(\P)$。見て分かるように、グラフ$G(\P)=(\{1,\dots,5,1',\dots,6'\},S(\P))$は森であり、木グラフの和集合として表せる。この場合は三つの木からなる。}
\end{figure}


$s$を$S(\P)$の濃度とする。$\P$が端点であるため、$s\leq n+m-1$である。$G(\P)$が巡回を持たないため、$G(\P)$は図~\ref{fig-f_of_P}に示されるように木または森（木の和集合）である。$\P$に相補的な対$(\fD,\gD)$を目指して、$n+m$個の変数に対する以下の$s$本の線形等式制約を考える：
\begin{equation}\label{eq-dual-variables}
	\begin{array}{ccc}
\fD_{i_1}+\gD_{j_1}&=&\C_{i_1,j_1}\\
\fD_{i_2}+\gD_{j_1}&=&\C_{i_2,j_1}\\
\vdots & = & \vdots\\
\fD_{i_s}+\gD_{j_s}&=&\C_{i_s,j_s},\\
\end{array}
\end{equation}
ここで$S(\P)$の要素は$(i_1,j_1'),\dots,(i_s,j_s')$と列挙されている。

$s\leq n+m-1 < n+m$であるから、上記の線形系~\eqref{eq-dual-variables}は常に不定である。この退化は、$\CouplingsD(\a,\b)$の$n+m$個の制約によるパラメータ化が$n+m$個の双対変数をもたらすことから部分的に解釈できる。注~\ref{rem-transportation-polytope}で概説されたより注意深い定式化では、$n+m-1$個の制約のみを持つ等価な定式化が得られ、したがって$n+m-1$個の双対変数となる。しかし、$s$は$n+m-1$よりも真に小さくなることもある：これは$G(\P)$が二つ以上の木の非連結和である場合に生じる。例えば、図~\ref{fig-f_of_P}では$5+6=11$個の双対変数（各ノードに一つ）があるが、これら$11$個のノード間の辺は$8$本、すなわち$(\fD,\gD)$を定義する線形方程式は$8$本しかない。したがって、この設定では$G(\P)$の連結成分の数だけ不定な双対変数が存在する。

$G(\P)$に列挙された木の一つを考える。その木がソースノード間に$k$個のノード$i_1,\dots,i_k$、ターゲットノード間に$l$個のノード$j_1',\dots,j_{l}'$を持ち、$r\eqdef k+l$ノードと$r-1$本の辺を持つとする。これは$\fD$の$k$個の変数と$\gD$の$l$個の変数に対応し、$r-1$本の線形方程式で結ばれている。不定性を解消するには、その木の根ノードを任意に選び、対応する双対変数に値$0$を割り当てればよい。そこから、幅優先探索または深さ優先探索を用いて木を走査し、その木内の他のすべての双対変数の値を決定する一連の単純な変数代入を得ることができる。これは図~\ref{fig-tree-linear-solve}に示されている。この手順を$\P$のグラフ内のすべての木に対して繰り返すことで、$\P$に相補的な双対変数の対$(\fD,\gD)$が得られる。

\begin{figure}[h!]
	\centering
	\includegraphics[width=\linewidth]{algo-basics-tree-solution/tree-linear-solve}
	\caption{\label{fig-tree-linear-solve}
 図~\ref{fig-f_of_P}に示されたグラフ$G(\P)$の最初の木に現れる五つのノードに対応する五つの双対変数$\fD_1,\fD_2,\gD_1,\gD_2,\gD_3$は、コスト行列$\C$の対応する要素を含む四つの線形方程式で結ばれている。この系は退化しているため、その木の根を選び（この例ではノード1）、対応する変数を$0$に設定し、根から木を走査（幅優先または深さ優先）することで、残りの四つの双対変数の値を反復的に求める。}
\end{figure}

\subsection{ネットワークシンプレックスの更新}
前節で得られた双対対$(\fD,\gD)$は、すべての$i,j$に対して$\fD_i+\gD_j\leq \C_{i,j}$が成り立つという意味で実行可能である可能性がある。その場合、命題~\ref{prop-primal-dual-optimality}により最適解に到達している。そうでない場合、すなわち$\fD_i+\gD_j> \C_{i,j}$となる$i,j$が存在する場合、ネットワークシンプレックスアルゴリズムが作動する。まず、グラフ$G$を実行可能解$\P$に対応するグラフ$G(\P)$に等しく初期化し、違反辺$(i,j')$を$G$に追加する。すると二つの場合が生じ得る：
\begin{enumerate}[label={(\alph*)}]
	\item $G$が（依然として）森である場合。これは$(i,j')$が既存の二つの部分木を繋ぐ場合に生じ得る。\S\ref{subsec-obtaining}で概説したアプローチをグラフ$G$に適用して新しい相補的双対ベクトル$(\fD,\gD)$を得ることができる。この追加は$n+m$個の双対変数の不定性を一つ除去するだけであり、主変数$\P$には何ら変更をもたらさないことに注意する。$(i,j')$はグラフ$G$に入ったが$\P_{i,j}$は$0$のままであるという意味で、この更新は通常退化と呼ばれる。ただし、$G(\P)$は$G$に含まれている。

	\item $G$が巡回を持つ場合。この場合、$G$が依然として森であることを保証するために$G$から辺を除去する必要があり、同時に$\P$を修正して$\P$が実行可能かつ$G(\P)$が$G$に含まれるようにする必要がある。これらの操作はすべて、$\P_{i,j}$の値を増加させ、検出された巡回に現れる$\P$の他の要素を修正することで実行できる。その方法は、命題~\ref{prop-extremal}の証明で用いたものと非常に類似している。より正確には、その巡回を$(i_1,j_1'), (j_1',i_2), (i_2,j_2'),\dots, (i_l,j_l'),(j_l',i_{l+1})$と書く。ただし、$i_1=i_{l+1}=i$として経路が$i$で始まり$i$で終わる巡回であることを保証し、$j_1=j$として巡回が追加された辺$\{i,j\}$から正の方向に始まることを強調する。すべての「正の」辺$(i_k,j_k')$（$k\leq l$に対して）のフローを増加させ、「負の」辺$(j_k',i_{k+1})$（$k\leq l$に対して）のフローを減少させて、更新された主解$\tilde{\P}$を得る。これは以下の要素を除いて$\P$に等しい：
$$\forall k\leq l,\quad \tilde{\P}_{i_k,j_k}:= \P_{i_k,j_k}+\theta;\quad \tilde{\P}_{i_{k+1},j_k}:= \P_{i_{k+1},j_k}-\theta.$$
ここで、$\theta$はその巡回を用いて添字$i,j$で可能な最大の増加量である。$\theta$の値は、巡回により負の影響を受ける最小のフロー、すなわち$\min_{k} \P_{i_{k+1},j_k}$で制御される。この更新は図~\ref{fig-simplex-update}に示されている。$k^\star$をその最小値を達成する添字とする。次に、$G$から$(i_{k^\star+1},j_{k^\star})$を除去し、\S\ref{subsec-obtaining}で概説したアプローチを用いて新しい双対変数$(\fD,\gD)$を計算することで更新を終了する。
\end{enumerate}

\subsection{主解の改善} 必ずしも当初の動機ではなかったが、上記の操作が$\P$のコストを改善することしかできないことを示すことができる。追加された辺が巡回を生じなかった場合、すなわち上記の場合{(a)}では、主解は変更されない。巡回が生じた場合、すなわち場合{(b)}では、$\P$は$\tilde{\P}$に更新され、以下の等式が成り立つ：
\eq{
\dotp{\tilde{\P}}{\C}-\dotp{\P}{\C}= \theta \left(\;\sum_{k=1}^{l} \C_{i_k,j_k} - \sum_{k=1}^{l}\C_{i_{k+1},j_k}\right).
}
ここで、前回の反復の最後に計算された双対ベクトル$(\fD,\gD)$を用いる。これらは$\f_{i_k}+\g_{i_k}=\C_{i_k,j_k}$かつ$\f_{i_{k+1}}+\g_{i_k}=\C_{i_{k+1},j_k}$をもともと$G$にあった\emph{すべての}辺に対して満たし、以下の等式をもたらす：
\eq{
\begin{aligned}\sum_{k=1}^{l} \C_{i_k,j_k} - \sum_{k=1}^{l}\C_{i_{k+1},j_k} &= \C_{i,j} + \sum_{k=2}^{l} \fD_{i_k}+\gD_{j_k} - \sum_{k=1}^{l} \fD_{i_{k+1}}+\gD_{j_k}\\
	&= \C_{i,j} - (\fD_{i} +\gD_{j}).
\end{aligned}
}
この項は、定義上、$i,j$が$C_{i,j} < \fD_{i} -\gD_{j}$であるために選ばれたことから負である。したがって、$\theta>0$ならば、\eq{\dotp{\tilde{\P}}{\C}=\dotp{\P}{\C}+\theta\left(\C_{i,j}-(\fD_{i} -\fD_{g})\right)<\dotp{\P}{\C}.}
$\theta=0$の場合（$G$と$G(\P)$が異なる場合に生じ得る）、グラフ$G$が変更されるだけで$\P$は変化しない。


\begin{figure}[h!]
	\centering
	\includegraphics[width=\linewidth]{algo-basics-tree-solution/simplex-update}
	\caption{\label{fig-simplex-update}
辺$\{i,j\}$をグラフ$G(\P)$に追加すると、(a) この追加後もグラフが森のままとなる場合があり、その場合は\S\ref{subsec-obtaining}で概説したアプローチに従って$\fD,\gD$を再計算できる；(b.1) この辺の追加により巡回が生じ、そこから有向経路を定義できる；(b.2) その経路を用いて$\P_{i,j}$の値を増加させ、フローの実行可能性制約を維持するためにその変更を巡回に沿って伝播させることができる。これは巡回により負の影響を受ける辺の一つのフローが$0$に減少するまで続く。これにより巡回が除去され$\P$が更新される。}
\end{figure}

したがって、ネットワークシンプレックスアルゴリズムは以下のように要約できる：\S\ref{subsec-northwest}で述べたNW隅法などにより与えられる端点解$\P$でアルゴリズムを初期化する。グラフ$G$を$G(\P)$で初期化する。\S\ref{subsec-obtaining}で述べた$G$内の木構造を用いた線形系の解法により、$\P$に相補的な双対変数の対$(\fD,\gD)$を計算する。{(i)} 制約$\C- \fD\oplus \gD\geq 0$に対する違反添字対を探す；見つからなければ$\P$は最適であり停止する。違反対$(i,j')$があれば、{(ii)} 辺$(i,j')$を$G$に追加する。$G$が依然として巡回を持たなければ、$(\fD,\gD)$をそれに応じて更新する；巡回があれば、$(i,j')$が正としてラベル付けされるように向きを付け、その巡回中の最小フロー値を持つ負の辺を除去し、図~\ref{fig-simplex-update}に示されるように$\P,G$を更新し、それに応じて相補的な対$\fD,\gD$を構築する；{(i)}に戻る。上記の操作の一部はグラフ操作（巡回検出、木走査）を必要とし、この文脈で効率的に実装できる。詳細は(\cite[\S5]{bertsekas1998network})に述べられている。

\citet{Orlin1997}がネットワークシンプレックスの多項式時間計算量を最初に証明した。\citet{Tarjan1997}はその直後に、ピボット辺の選択を支援するより効率的なデータ構造に依拠する改良された計算量$O\left(\,(n+m)nm \log(n+m)\log\left((n+m)\|\C\|_\infty\right)\,\right)$を示した。



%\paragraph{Improvement on $\P$.} Introducing a new edge $\{i,j\}$ in $G(\P)$ only because that pair violated the complementary requirement, namely because $\fD_i+\gD_j$ is strictly larger that $\C_{i,j}$, may seem like a valid approach to get closer to the goal of obtaining a feasible pair $(\fD,\gD)$ but it does not necessarily explain why this would result in a better solution $\P$ at all.

\section{双対上昇法}\label{s-dual-ascent}

双対上昇法はネットワークシンプレックスに数十年先立つものであり、\citet{borchardt1865investigando}の研究、および後に\citet{Kuhn1955}によって再述されたK\"onigとEgerv\'aryの研究にまで遡ることができる。ハンガリアン法はこの族で最もよく知られたアルゴリズムであり、$\a$と$\b$が等しくともに一様分布である特殊な場合、すなわち$\a=\b=\ones_n/n$の場合にのみ動作する。以下では、より一般的な双対上昇法の族の簡潔な説明を行う。これには最大流問題の知識が必要である(\cite[\S7.5]{bertsimas1997introduction})。上述の主問題で提示されたネットワークシンプレックスとは対照的に、双対上昇法は各反復において双対実行可能解を維持し、$\fD$と$\gD$に疎なベクトルを加えることでその目的関数値を徐々に改善する。我々の説明は主に(\cite[\S7.7]{bertsimas1997introduction})に基づいており、以下の定義から始まる。

\begin{defn} $S\subset\range{n},S'\subset\range{m}'\eqdef \{1',\dots,m'\}$に対して、$\ones_S$を$S$に列挙される添字で$1$、それ以外で$0$の$\RR^n$のベクトル、同様に$\ones_{S'}$を$S'$の添字を持つ$\RR^m$のベクトルと書く。
\end{defn}

以下では、$(\fD,\gD)$を$\PotentialsD(\C)$の実行可能な双対対とする。これは、すべての対$(i,j')\in\range{n}\times\range{m}'$に対して$\fD_i+\gD_j\leq\C_{ij}$が成り立つことを意味する。$(i,j')$が$\fD_i+\gD_j=\C_{ij}$のとき\emph{均衡}対（または辺）と呼び、$\fD_i+\gD_j<\C_{ij}$のとき\emph{非活性}と呼ぶ。この規約のもと、実行可能な双対対$(\fD,\gD)$が集合$S$と$S'$で添字付けられた疎なベクトルを用いて摂動されても実行可能であり続ける条件を記述する簡単な結果から始める。

\begin{prop}\label{prop-feasibility-dual-pair}
$(\tilde{\fD},\tilde{\gD})\eqdef(\fD,\gD)+\varepsilon(\ones_S,-\ones_{S'})$が十分小さい$\varepsilon>0$に対して双対実行可能であるための条件は、すべての$i\in S$に対して$(i,j')$が均衡であれば$j'\in S'$であることである。
\end{prop}
\begin{proof}
任意の$i\in S$に対して、$(i,j')$が非活性であるすべての$j'\in\range{m}'$、すなわち$\fD_i+\gD_j<\C_{ij}$であるすべての$j'$の集合$\mathcal{I}_{i}$を考える。$\varepsilon_i \eqdef \min_{j\in I_{i}} \C_{i,j}-\fD_i-\gD_j$を、$j'\in \mathcal{I}_i$に対応する制約を違反することなく$\fD_i$を増加させることのできる最小のマージンと定義する。実際、$\varepsilon\leq \varepsilon_i$ならば任意の$j'\in \mathcal{I}_i$に対して$\tilde{\fD}_i+\tilde{\gD}_j<\C_{i,j}$が成り立つ。次に、$i$に関連する均衡辺の集合$\mathcal{B}_i$を考える。$\mathcal{B}_i=\range{m}'\setminus \mathcal{I}_i$であることに注意する。上記の仮定は$j'\in\mathcal{B}_i\Rightarrow j'\in S'$である。したがって、$j'\in\mathcal{B}_i$に対して$\tilde{\fD}_i+\tilde{\gD}_j=\fD_i+\gD_j=\C_{i,j}$が成り立つ。結果として、不等式$\tilde{\fD}_i+\tilde{\gD}_j\leq \C_{i,j}$は任意の$j\in \range{m}'$に対して保証される。許容される最小の増加量、すなわち$\min_{i\in S} \varepsilon_i$よりも小さい増加$\varepsilon$を選ぶと、$(\tilde{\fD},\tilde{\gD})$が双対実行可能であることが回復される。
\end{proof}

\S\ref{subsec-obtaining}で示したネットワークシンプレックスの反復の主な動機は、実行可能な主解$\P$から出発して相補的な実行可能双対対$(\fD,\gD)$を得ることである。その目標を達成するため、$\P$はその相補的双対対が双対実行可能性に到達するよう徐々に修正される。対称的なアプローチ、すなわち実行可能な双対変数から出発して実行可能な主変数$\P$を得るアプローチが、双対上昇法の動機となっている。以下の命題は双対上昇法の主要なエンジンであり、$(\fD,\gD)$に対して実行可能性を維持する上昇方向の存在を（構成的に）保証する。その方向は、ネットワークシンプレックスと同様に、候補の主解$\P$を設計し、その非実行可能性が$(\fD,\gD)$の更新を導くことで構築される。

\begin{prop}\label{prop-dual-ascent-dir} $(\fD,\gD)$が問題~\eqref{eq-dual-discrete-split}に対して最適であるか、または$S\subset\range{n},S'\subset\range{m}'$が存在して$(\tilde{\fD},\tilde{\gD})\eqdef(\fD,\gD)+\varepsilon(\ones_S,-\ones_{S'})$が十分小さい$\varepsilon>0$に対して実行可能でありかつ真に良い目的関数値を持つかのいずれかである。
\end{prop}
\begin{proof}
まず、$(\fD,\gD)$に相補的な主変数$\P$を考える。そのために、均衡辺の集合$\mathcal{B}$、すなわち$\fD_i+\gD_j=\C_{i,j}$であるすべての対$(i,j')\in\range{n}\times\range{m}'$を定め、頂点$\{1,\dots,n,1',\dots,m'\}$が$\mathcal{B}$の辺のみで結ばれた二部グラフを形成する。これに、ソースノード$s$をすべてのノード$i\in\range{n}$に\emph{容量制限付き}辺（容量$\a_i$）で接続し、ターミナルノード$t$をすべてのノード$j'\in\range{m}'$に容量$\b_j$の辺で接続して補完する。図~\ref{fig-dualascent}に示される通りである。Ford--Fulkersonアルゴリズム(\cite[p. 305]{bertsimas1997introduction})をこのネットワークに適用して最大流$\mathbf{F}$を計算できる。すなわち、$(i,j')\in \mathcal{B}$で添字付けられる$n+m+|\mathcal{B}|$個の非負値の族$f_{si}\leq \a_i, f_{ij'}, f_{j't}\leq\b_j$で、流量保存制約を満たし$\sum_i f_{si}$が最大であるものを求める。この流$\mathbf{F}$のスループットが$1$に等しければ、$\fD,\gD$に構成上相補的な実行可能主解$\P$が$\mathbf{F}$から抽出でき、$(i,j')\in \mathcal{B}$に対して$\P_{i,j}=f_{ij'}$、それ以外では零と定義する。これにより命題~\ref{prop-primal-dual-optimality}から$(\fD,\gD)$と$\P$の最適性が得られる。$\mathbf{F}$のスループットが$1$より真に小さい場合、ラベリングアルゴリズムは$s$から反復的に到達可能で$\mathbf{F}$が容量制約を飽和させないノード、およびラベル付けされたノードにフローを供給するノードをラベル付け（識別）する。ラベル付けされたノードは空でない集合$Q$に格納され、$\mathbf{F}$の最適性によりターミナルノード$t$は含まれない（厳密な提示は\citealt[p. 308]{bertsimas1997introduction}を参照）。$Q$は二つの集合$S=Q\cap\range{n}$と$S'=Q\cap\range{m}'$に分割できる。総スループットが$1$より真に小さいと仮定しているため、$S\ne\emptyset$である。まず、$i\in S$かつ$(i,j)$が均衡であれば、$j'$は必然的に$S'$に属することに注意する。実際、すべての辺$(i,j')$は構成上無限容量を持つため、ラベリングアルゴリズムは$i$を$S$に含めると必然的に$j'$にも到達する。命題~\ref{prop-feasibility-dual-pair}により、$\tilde{\fD},\tilde{\gD}$の実行可能性を保証する十分小さい$\varepsilon$が存在する。$(\tilde{\fD},\tilde{\gD})$が$(\fD,\gD)$よりも良い目的関数値を持つことを保証するためには、$\transp{\ones_S}{\a}-\transp{\ones_{S'}}{\b}>0$を証明する必要がある。$\bar{S}=\range{n}\setminus S$、$\bar{S'}=\range{m}'\setminus S'$とし、
$$A = \sum_{i\in S} f_{si}, \quad B = \sum_{i\in \bar{S}} f_{si}, \quad C = \sum_{j'\in S'} f_{j't}, \quad D = \sum_{j'\in \bar{S'}} f_{j't}.$$
と定義する。総最大流は$s$から出発するため$A+B$に等しく、同時に$t$に到着するため$C+D$にも等しい。流量保存制約により、同じ流はまた$B+C$にも等しく、したがって$A=C$である。一方、ラベリングアルゴリズムの定義により、$S$のすべての$i$に対して$f_{si}<\a_i$であり、$\bar{S}'$の$j'$に対しては$t$が最適流により$S'$に含まれないため$f_{j't}=\b_j$である。したがって$A<\transp{\ones_S}\a$かつ$C=\transp{\ones_S'}\b$である。よって$\transp{\ones_S}\a-\transp{\ones_S'}\b>A-C=0$が成り立つ。
\end{proof}

\begin{figure}[h!]
\centering
\includegraphics[width=\linewidth]{algo-basics-tree-solution/dual_ascent}
\caption{\label{fig-dualascent}
図~\ref{fig-f_of_P}で最初に導入した周辺分布を含む輸送問題を考える。$n=5, m=6$である。二つの実行可能な双対ベクトル$\fD,\gD$が与えられたとき、$(\fD,\gD)$に相補的な「最良の」フロー行列$P$を求めようとする。これは$\P$が$\fD_i+\gD_j=\C_{i,j}$となる添字に対応する辺$(i,j')$でのみ正の値を取れることを意味する。これはプロット(a)で点線で表されている。そのグラフ構造で達成できる最良のフローは、抽象的なソースノード$s$から$i\in\range{n}$とラベル付けされたすべてのノードへの辺、抽象的なターミナルノード$t$から$j\in\range{m'}$の$j'$とラベル付けされたすべてのノードへの辺を持つ容量制限付きネットワークにおける最大流問題として定式化できる。ここで辺$(s,i),(j',t), i\in \range{n},j\in\range{m}$の容量はそれぞれ$\a_i,\b_j$であり、他はすべて無限である。Ford--Fulkersonアルゴリズム(\cite[p. 305]{bertsimas1997introduction})を適用してそのような最大流を計算できるが、プロット(b)に示されるように、問題を解くために必要な$1$単位の質量のうち$0.74$単位しか達成できない。最大流アルゴリズムで使われるサブルーチンの一つであるラベリングアルゴリズム(\cite[p. 308]{bertsimas1997introduction})を用いて、$s$から飽和していないフローを受け取るノード（および再帰的にそのすべての後続ノード）を識別でき、プロット(c)ではオレンジの線で示されている。ラベリングアルゴリズムはまた、ラベル付けされたノードに正のフローを送るノードもデフォルトで追加する。この基準によりノード$3$が選択され、$3'$への赤い線で示されている。ラベル付けされたノードは集合$S,S'$にグループ化でき、$\fD,\gD$を変更して異なるグラフを得ることで、より高いフローを得るためにより良く活用できるノードを識別する。証明にはプロット(d)に記述された部分フロー和が関わる。}
\end{figure}

双対上昇法は、実行可能性を保証し目的関数を改善する集合$S,S'$により生成される任意のベクトルで実行可能解$(\fD,\gD)$を修正することで進む。集合$S,S'$が命題~\ref{prop-dual-ascent-dir}の証明で構成的に与えられたものであり、ステップ長$\varepsilon$が命題~\ref{prop-feasibility-dual-pair}の証明で定義されたものである場合、\emph{主双対法}として知られる方法が回復される。この方法はマッチング問題に対してはハンガリアン法に帰着される。双対上昇法はネットワークシンプレックスの双対変形と類似点を持つが、少なくとも二つの重要な点で異なる。シンプレックス型の方法は、現在の解が常に実行可能集合の\emph{端点}（双対の場合は$\PotentialsD(\C)$）であることを保証するが、ここで提示された双対上昇法はそのような仮定を置かず、実行可能集合の内部にある反復点を自由に生成できる。さらに、双対ネットワークシンプレックスが$(\fD,\gD)$を修正して線形（周辺制約）を満たす主解$\P$を生成し、非負性は収束時にのみ満たされるのに対して、双対上昇法は常に非負であるが必ずしも周辺制約を満たさない主解$\P$を構築する。

\section{競売アルゴリズム}\label{s-auction}
競売アルゴリズムは~\citet{bertsekas1981new}によって最初に提案され、後に~\citep{bertsekas1988dual}で改良された。このアルゴリズムにはいくつかの経済学的解釈が提案されている（例えば\citet{bertsekas1992auction}を参照）。このアルゴリズムは任意の周辺分布に対して適用可能であるが、ここでは最適割当問題を解くための定式化で提示する。

\paragraph{相補スラック性。} 最適割当問題では、最適輸送問題に対して示した主双対条件がより定式化しやすくなることに注意する。任意の端点解$\P$は\todoK{there always exist a permutation solution, but the converse it not true}必然的にある$\si$に対する置換行列$\P_\si$である（式~\eqref{prop-primal-dual-optimality}を参照）。最適な主変数$\P_{\si^\star}$と双対変数$\fD^\star,\gD^\star$が与えられると、
$$
\fD^\star_i + \gD^\star_{\si^\star_i} = \C_{i,{\si_i^\star}}.
$$
が必然的に成り立つ。また、\S\ref{sec-c-transforms}で述べた$\C$変換の原理により、$\fD^\star$は$\gD^{\bar\C}$に等しく選ぶことができることを想起されたい。したがって、
\eql{\label{eq-dual-conditions-assign-1}
\C_{i,\si^\star_i} - \gD^\star_{\si_i}  = \min_{j} \C_{i,j} - \gD^\star_{j}.
}
逆に、あるベクトル$\gD$と置換$\sigma$が存在して
\eql{\label{eq-dual-conditions-assign}
\C_{i,\si_i} - \gD_{\si_i}  = \min_{j} \C_{i,j} - \gD_{j}
}
が成り立つならば、それらはともに最適であること、すなわち$\si$が最適割当であり$\gD^{\bar\C},\gD$が最適双対対であることが容易に示される。


\paragraph{部分割当と$\varepsilon$-相補スラック性。} 競売アルゴリズムの目標は、三つ組$S,\xi,\gD$を反復的に修正することである。ここで$S$は$\range{n}$の部分集合、$\xi$は部分割当ベクトル（$S$から$\range{n}$への単射）、$\gD$は双対ベクトルである。双対ベクトルは\emph{近似的な}相補スラック性~\eqref{eq-dual-conditions-assign}を満たす解に収束することを意図しており、$S$は$\xi$が置換を記述するにつれて$\range{n}$を覆うように成長する。アルゴリズムは各反復後に以下の三つの性質を維持することで動作する：
\begin{enumerate}[label={(\alph*)}]
	\item $\forall i\in S,\quad \C_{i,\xi_i} - \gD_{\xi_i}  \leq \varepsilon + \min_{j} \C_{i,j} - \gD_{j}$（$\varepsilon$-CS）。
	\item $S$のサイズは各反復で増加することしかできない。
	\item ある添字$i$が存在して$\gD_i$が少なくとも$\varepsilon$だけ減少する。
\end{enumerate}

\paragraph{競売アルゴリズムの更新。} 点$j$が与えられると、競売アルゴリズムは通常の$\C$変換に現れる最適値だけでなく、二番目に良い値も使用する：
$$j^1_i \in \argmin_{j} \C_{i,j} - \gD_{j}, \quad j^2_i \in \argmin_{j\ne j^1_i} \C_{i,j} - \gD_{j}, $$
以下の更新を、$i\notin S$である添字$i$に対する$\gD$の更新、および$S$と$\xi$の更新として定義する：
\begin{enumerate}\item \textbf{$\gD$の更新}：$\gD$の第$j^1_i$要素から、$\varepsilon$と調整コスト$\{\C_{i,j}-\gD_j\}_j$の最低値と二番目の最低値の差の和を引く：
\eql{
\begin{aligned}
	\gD_{j^1_{i}} &\leftarrow \gD_{j^1_{i}} - \underbrace{\left((\C_{i,j^2_{i}}-\gD_{j^2_{i}})-(\C_{i,j^1_{i}}-\gD_{j^1_{i}})+\varepsilon\right)}_{\geq \varepsilon >0}\label{eq-auction-update}\\
	&= \C_{i,j^1_{i}} - (\C_{i,j^2_{i}}-\gD_{j^2_{i}}) -\varepsilon.\\
\end{aligned}
}
\item \textbf{$S$と$\xi$の更新}：$\xi_{i'}=j^1_{i}$となる添字$i'\in S$が存在すれば、$S\leftarrow S\setminus \{i'\}$として除去する。$\xi_{i}=j^1_{i}$と設定し、$i$を$S$に追加する。$S\leftarrow S \cup \{i\}.$
\end{enumerate}

\paragraph{アルゴリズムの性質。} アルゴリズムは割当点の空集合$S=\emptyset$、割当なし、空の部分割当ベクトル$\xi$、および$\gD=\zeros_n$から開始し、$S=\range{n}$で終了し、終了するまで上記の両ステップを繰り返す。性質{(b)}と{(c)}が各反復後に成り立つことは、更新の性質から明らかである（式~\eqref{eq-auction-update}を見れば十分である）。$\varepsilon$-相補スラック性は最初の反復で容易に満たされる。この場合$S=\emptyset$であるためである。反復がこの性質を保存するという事実は以下の命題で示される。
%Show::: $(\fD^{\,\tau})_i < (\fD^{\,\C\bar{\C}})_i$
\begin{prop} 競売アルゴリズムは各反復において$\varepsilon$-相補スラック性を維持する。
\end{prop}
\begin{proof}
	$\gD,\xi,S$をある反復の開始時の三つの変数とする。したがって、任意の$i'\in S$に対して
	$$\C_{i,\xi_{i'}} - \gD_{\xi_{i'}}  \leq \varepsilon + \min_{j} \C_{{i'},j} - \gD_{j} $$
が成り立つと仮定する。ここで反復で考慮される特定の$i\notin S$を考える。三つの更新が行われる：$\gD,\xi,S$は添字$j^1_i$と$j^2_i$を用いて$\gD^{\text{n}},\xi^{\text{n}},S^{\text{n}}$に更新される。より正確には、$\gD^{\text{n}}$は要素$j^1_i$を除いて$\gD$に等しく、その値は
$$\gD_{j^1_{i}}^\text{n} = \gD_{j^1_{i}} - \left((\C_{i,j^2_{i}}-\gD_{j^2_{i}})-(\C_{i,j^1_{i}}-\gD_{j^1_{i}})\right)-\varepsilon \leq \gD_{j^1_{i}}-\varepsilon$$,
$\xi^{\text{n}}$は第$i$要素が$j^1_i$に等しい点を除いて$\xi$に等しく、$S^{\text{n}}$は$\{i\}$と$S$の和集合に等しい（一つの要素が除去される可能性がある）。$\gD^{\text{n}}$の更新は次のように書き直せる：
$$\gD^{\text{n}}_{j^1_i}= \C_{i,j^1_{i}} - (\C_{i,j^2_{i}}-\gD_{j^2_{i}}) -\varepsilon;$$
したがって
$$\C_{i,j^1_{i}} - \gD^{\text{n}}_{j^1_i}=  \varepsilon + (\C_{i,j^2_{i}}-\gD_{j^2_{i}}) = \varepsilon+\min_{j\ne j^1_i} (\C_{i,j}-\gD_{j}).$$
$-\gD\leq -\gD^{\text{n}}$であるから、
$$\C_{i,j^1_{i}} - \gD^{\text{n}}_{j^1_i} = \varepsilon + \min_{j\ne j^1_i} (\C_{i,j}-\gD_{j})\leq \varepsilon+\min_{j\ne j^1_i} (\C_{i,j}-\gD^{\text{n}}_{j}),$$
が得られ、$j=j^1_i$に対しても不等式は明らかに成り立つため、添字$i$に対する$\varepsilon$-相補スラック性が得られる。他の添字$i'\ne i$に対しては、同様に$\gD^{\text{n}}\leq\gD$であるから、以下の不等式の列が成り立つ：
$$\C_{i,\xi^{\text{n}}_{i'}} - \gD^{\text{n}}_{\xi^{\text{n}}_{i'}} = \C_{i,\xi_{i'}} - \gD_{\xi_{i'}} \leq \varepsilon + \min_{j} \C_{{i'},j} - \gD_{j} \leq \varepsilon + \min_{j} \C_{{i'},j} - \gD^{n}_{j}. $$
\end{proof}

\begin{prop} 競売アルゴリズムのステップ数は最大$N=n\|\C\|_\infty/\varepsilon$である。
\end{prop}
\begin{proof}アルゴリズムが$T>N$ステップ後に停止していないと仮定する。すると、$\xi$の像に含まれない添字$j$が存在し、その価格座標$\gD_j$は一度も更新されず$\gD_j=0$のままである。この場合、$n>\|\C\|_\infty/\varepsilon$回更新された添字$j'$は存在し得ない。実際、もしそうであれば任意の添字$i$に対して
$$\gD_{j'}\leq-n\varepsilon<-\|\C\|_\infty\leq -\C_{i,j} = \gD_j-\C_{i,j},$$
が成り立ち、すべての$i$に対して
$$\C_{i,j'}-\gD_{j'}> \C_{i,j}+ (\C_{i,j}-\gD_j),$$
となるが、これは$\varepsilon$-CSに矛盾する。したがって、各変数の更新回数は$\|C\|_\infty/\varepsilon$を超えないため、総反復回数$T$は$n\|\C\|_\infty/\varepsilon=N$を超えることができない。
\end{proof}

\begin{rem} この結果は、アルゴリズムが終了するための素朴な演算回数として$N^3\|\C\|_\infty/\varepsilon$を与えることに注意されたい。この計算量は、$\varepsilon$の値を各反復で減少させる巧妙な方法である$\varepsilon$スケーリングを用いることで$N^3 \log\|\C\|_\infty$に削減できる(\cite[p. 264]{bertsekas1998network})。
\end{rem}

\begin{prop} 競売アルゴリズムは、コストが$n\varepsilon$だけ準最適な割当を見つける。
\end{prop}
\begin{proof}
	$\sigma,\gD^\star$をコスト行列$\C$の割当問題の主および双対最適解とし、最適値を$$t^\star=\sum \C_{i,\sigma_i} = \sum_i \min_{j}\C_{i,j}-\gD^\star_j+\sum_j\gD^\star_j.$$とする。
$\xi,\gD$を終了時に競売アルゴリズムが出力する解とする。$\varepsilon$-CS条件により、任意の$i\in S$に対して
$$\min_{j} \C_{i,j} - \gD_{j}\geq \C_{i,\xi_i} - \gD_{\xi_i} -\varepsilon.$$
が成り立つ。したがって、$\gD$の準最適性によりまず
	$$\begin{aligned}t^\star&\geq \sum_i \left(\min_{j} \C_{i,j}-\gD_{j}\right) + \sum_j \gD_j\\
&\geq
	\sum_i -\varepsilon+\left(\C_{i,\xi_i}-\gD_{\xi_i}\right) + \sum_j \gD_j = -n\varepsilon + \sum_i \C_{i,\xi_j} \geq  - n\varepsilon + t^\star.\end{aligned},$$
が得られる。ここで第二の不等式は$\varepsilon$-CSから、次の等式は$\gD_{\xi_i}$の項と$\gD_{j}$の和の打消しから、最後の不等式は置換としての$\xi$の準最適性から従う。
\end{proof}


したがって、競売アルゴリズムは$\C$変換の機構を利用する別の方法と見なすことができる。次に、正則化に基づく別のアプローチ、いわゆるSinkhornアルゴリズムを探究する。これは~\citep{schmitzer2016stabilized}で議論されているように、競売アルゴリズムとも類似点がある。
% $\fD^{\,\C\bar{\C}\C}\geq \fD^{\,\C}$ implies that $\fD^{\,\C\bar{\C}\C\bar{\C}}\leq \fD^{\,\C\bar{\C}}$
%

最後に、Euclid空間の低次元正則格子上では、これらの古典的な線形ソルバーにマルチスケール戦略を組み合わせることで、大幅な高速化が可能であることを注記する~\citep{schmitzer2016sparse,oberman2015efficient}。




%\section{Interior point methods}\label{s-ipm}
% \begin{rem}\label{rem:duality}Recall that a simple argument can help us recover the dual problem from the primal. For convenience, let us write $\q=\bigl[\begin{smallmatrix}\a\\ \b \end{smallmatrix} \bigr]$.
% 	Let $\U$ be the feasible set $\{\p\in\mathbb{R}^{nm}, \mathbf{A}\p = \q\}$.
% 	The feasible set's indicator function $\iota_\U$, such that $\iota_{\U}(\p)=0$ if $\p\in \U$ and $\infty$ otherwise,
% 	can be written as $\max_{\hD\in\mathbb{R}^{n+m}} -\hD^T(\mathbf{A}\p-\q)$.
% 	Indeed, the optimal value of that problem is obviously 0 when $\p$ is such that $\mathbf{A}\p=\q$, and $\infty$ otherwise, since it suffices to take for $\hD$ any negative constant that is large enough, and multiply it by vector $\mathbf{A}\p-\q$ to obtain an arbitrarily large objective value. Therefore, using the extended definition of convex functions, we have that
% $$\begin{aligned}
% \umin{\substack{\p \in \RR^{nm}_+\\\mathbf{A}\p = \q}} \cc^T\p & = \umin{\p \in \RR^{nm}_+} \cc^T\p + \iota_S(\p) =\umin{\p \in \RR^{nm}_+} \cc^T\p + \max_{\hD\in\mathbb{R}^{n+m}} -\hD^T(\mathbf{A}\p-\q)
% \\ & = \umin{\p \in \RR^{nm}_+}  \max_{\hD\in\mathbb{R}^{n+m}}  \cc^T\p - \hD^T(\mathbf{A}\p-\q)
% \\ & = \max_{\hD\in\mathbb{R}^{n+m}} \hD^T\q + \umin{\p \in \RR^{nm}_+} (\cc-\transp{\mathbf{A}}\hD)^T\p.
% \end{aligned}$$
% where we have assumed in the third line that one can inverse the orders of $\min$ and $\max$ problems. One can also notice that the term $\umin{\p \in \RR^{nm}_+} (\cc-\transp{\mathbf{A}}\hD)^T\p$ is minus the indicator function of the set $\R=\{\hD, \transp{\mathbf{A}}\hD\leq \cc\}$. Indeed, if any element of vector $\transp{\mathbf{A}}\hD$ is strictly bigger than $\cc$, say at index $i$, then the $i$-th element of $\cc-\transp{\mathbf{A}}\hD$ is negative, and it suffices to take for $\p$ the null vector except for its $i$-th index with an arbitrary large value to obtain and arbitrarily small value for the minimum. Therefore we have that
% $$
% \umin{\substack{\p \in \RR^{nm}_+\\\mathbf{A}\p = \q}} \cc^T\p = \max_{\hD\in\mathbb{R}^{n+m}} (\transp{\mathbf{A}}\hD)^T\q -\iota_{\R}(\hD)= \max_{\substack{\hD\in\mathbb{R}^{n+m}\\ \transp{\mathbf{A}}\hD\leq \cc}}\hD^T\q
% $$
% For a thorough and rigorous duality argument we refer the reader to~\cite[\S4.3]{bertsimas1997introduction}:
% \end{rem}
